\scanpage{72}
and in terms of the $f\ldots$ the formula, is. identical but a little less friendly in appearance
\[
(11)\qquad f_{ll'}(u)=l'^{-1}\bigl(f_l(u)\bigr)\cdot f_{l'}(u)
\qquad\text{i.e. }f_u(ll')=l'^{-1}f_u(l)\cdot f_u(l')
\]

\deleted{Formula (10) can be expressed by saying that for $l$ fixed, $u\in\Gamma$ variable,
\[
\unsure{(17)|(12)}\qquad g_l:\Gamma\longrightarrow\widehat{\pi}_1(X,y)
\]
is a $1$-cocycle on $\Gamma$, with values in $\widehat{\pi}_1(X,y)$, relative to the natural action of $\Gamma$ on $\widehat{\pi}_1(X,y)$. Of course, the left $\widehat{\pi}_1(X,y)$-torsor, \inserted{with an action of $\Gamma$,} corresponding to this $1$-cocycle, is none other than $\operatorname{Hom}_{\widehat{\Pi}_1(X)}(x,y)$, on which $\widehat{\pi}_1(X,y)=\operatorname{Hom}_{\widehat{\Pi}_1(X)}(y,y)$ acts by composition.}

For a composite of a \deleted{finite number} finite composable sequence of paths
\diagram{\begin{tikzcd}[ampersand replacement=\&]
x_0\arrow[r,"l_1"]\&x_1\arrow[r,"l_2"]\&\cdots\&x_{n-1}\arrow[r,"l_n"]\&x_n
\end{tikzcd}}
formula (10) takes the form
\[
\unsure{(12)|(13)}\qquad
 g_{l_nl_{n-1}\cdots l_1}(u)=g_{l_n}(u)\cdot l_n\bigl(g_{l_{n-1}}(u)\bigr)\cdot l_nl_{n-1}\bigl(g_{l_{n-2}}(u)\bigr)\ldots
 \inserted{\cdot(l_n\cdots l_2)g_{l_1}(u))}
\]
The most \deleted{obvious} \inserted{simple} formula possible, in short, of those one could imagine that makes sense. We are going to develop \deleted{further} \inserted{a} context that allows one \deleted{better} to write this formula, as well as \deleted{(11)} the analogous formula in $f$ and the formulas
\margin{71}
